\documentclass[a4paper,11pt]{scrartcl}
% Преамбула раздатки мастер-класса. Сборка: latexmk -xelatex masterclass.tex
% Оформление то же, что у учебника (docs/uchebnik/preamble.tex).
\usepackage{fontspec}
\usepackage{polyglossia}
\setmainlanguage{russian}
\setotherlanguage{english}
\setmainfont{Charis SIL}[Scale=0.93]
\setsansfont{Fira Sans}
\setmonofont{JetBrains Mono}[Scale=0.84, RawFeature={-calt;-liga}]
\newfontfamily\cyrillicfonttt{JetBrains Mono}[Scale=0.84, RawFeature={-calt;-liga}]

\usepackage{xcolor}
\definecolor{accent}{HTML}{0B7285}
\definecolor{accentdark}{HTML}{095C6B}
\definecolor{accentsoft}{HTML}{E3F1F4}
\definecolor{ink}{HTML}{1F2933}
\definecolor{inkgray}{HTML}{5B6770}
\definecolor{rule}{HTML}{C9D2D6}
\definecolor{paper}{HTML}{F5F7F8}
\definecolor{warnfill}{HTML}{FBF3E4}
\definecolor{warnedge}{HTML}{B7791F}
\definecolor{stopfill}{HTML}{FAECEC}
\definecolor{stopedge}{HTML}{A33A3A}
\definecolor{ledred}{HTML}{D64545}
\definecolor{ledgreen}{HTML}{3C9D5A}
\definecolor{ledblue}{HTML}{3465C4}
\definecolor{ledyellow}{HTML}{D9B22B}
\definecolor{ledcyan}{HTML}{2AA0B8}

\usepackage{geometry}
\geometry{a4paper, top=22mm, bottom=24mm, left=22mm, right=22mm, headsep=8mm}
\KOMAoptions{parskip=half-}
\linespread{1.06}

\usepackage{microtype}
\usepackage{booktabs}
\usepackage{array}
\usepackage{tabularx}
\usepackage{enumitem}
\setlist{itemsep=2pt, topsep=4pt, parsep=0pt}
\setlist[enumerate,1]{leftmargin=1.9em}
\setlist[itemize,1]{leftmargin=1.3em}
\usepackage{caption}
\captionsetup{font={small,color=inkgray}, labelfont={bf,color=ink}, labelsep=period, justification=raggedright, singlelinecheck=false, skip=6pt}
\usepackage{fvextra}
\usepackage[most]{tcolorbox}
\usepackage{tikz}
\usetikzlibrary{arrows.meta, calc, positioning}
\usepackage{float}
\usepackage{needspace}
\raggedbottom

% ---------- заголовки
\addtokomafont{disposition}{\sffamily\color{ink}}
\setkomafont{section}{\sffamily\bfseries\Large\color{accentdark}}
\setkomafont{subsection}{\sffamily\bfseries\large\color{ink}}
\RedeclareSectionCommand[beforeskip=1.4\baselineskip plus 4pt, afterskip=.5\baselineskip]{section}
\RedeclareSectionCommand[beforeskip=1.0\baselineskip plus 3pt, afterskip=.3\baselineskip]{subsection}

\usepackage[headsepline]{scrlayer-scrpage}
\pagestyle{scrheadings}
\clearpairofpagestyles
\ihead{Мастер-класс: шоу из четырёх дронов}
\ohead{\pagemark}
\setkomafont{pagehead}{\sffamily\small\color{inkgray}}
\setkomafont{pagenumber}{\sffamily\small\bfseries\color{ink}}
\setkomafont{headsepline}{\color{rule}}

\usepackage[unicode, pdftitle={Мастер-класс: шоу из четырёх дронов}, pdfauthor={Litebee IDE}]{hyperref}
\hypersetup{colorlinks=true, linkcolor=accent, urlcolor=accent}

% ---------- листинги
\fvset{fontsize=\small, breaklines=true, numbers=left, numbersep=7pt, xleftmargin=16pt, tabsize=4}
\renewcommand{\theFancyVerbLine}{\textcolor{inkgray}{\scriptsize\arabic{FancyVerbLine}}}

\tcbset{
  codebox/.style={enhanced, breakable, colback=paper, colframe=paper, boxrule=0pt, sharp corners, frame hidden,
    borderline west={2.5pt}{0pt}{accent}, left=9pt, right=6pt, top=3pt, bottom=3pt,
    before skip=4pt, after skip=10pt},
  note/.style={enhanced, breakable, colback=accentsoft, colframe=accentsoft, frame hidden, sharp corners,
    borderline west={3pt}{0pt}{accent}, left=9pt, right=8pt, top=5pt, bottom=5pt, before skip=8pt, after skip=10pt},
  warn/.style={enhanced, breakable, colback=warnfill, colframe=warnfill, frame hidden, sharp corners,
    borderline west={3pt}{0pt}{warnedge}, left=9pt, right=8pt, top=5pt, bottom=5pt, before skip=8pt, after skip=10pt},
  stop/.style={enhanced, breakable, colback=stopfill, colframe=stopfill, frame hidden, sharp corners,
    borderline west={3pt}{0pt}{stopedge}, left=9pt, right=8pt, top=5pt, bottom=5pt, before skip=8pt, after skip=10pt},
  task/.style={enhanced, breakable, colback=white, colframe=accent, boxrule=0.9pt, arc=3pt,
    left=10pt, right=10pt, top=6pt, bottom=6pt, before skip=12pt, after skip=10pt,
    fonttitle=\sffamily\bfseries\normalsize, coltitle=white, colbacktitle=accent, toptitle=3pt, bottomtitle=3pt},
}

\newenvironment{kod*}{\VerbatimEnvironment\begin{tcolorbox}[codebox]\begin{Verbatim}}{\end{Verbatim}\end{tcolorbox}}

\newenvironment{zametka}[1][Запомните]{\begin{tcolorbox}[note]{\sffamily\bfseries\small\color{accentdark}#1}\par\smallskip}{\end{tcolorbox}}
\newenvironment{vnimanie}[1][Осторожно]{\begin{tcolorbox}[warn]{\sffamily\bfseries\small\color{warnedge}#1}\par\smallskip}{\end{tcolorbox}}
\newenvironment{opasno}[1][Безопасность]{\begin{tcolorbox}[stop]{\sffamily\bfseries\small\color{stopedge}#1}\par\smallskip}{\end{tcolorbox}}

\newcounter{zadanie}
\newenvironment{zadanie}[1]{\refstepcounter{zadanie}\begin{tcolorbox}[task, title={Задание \thezadanie. #1}]}{\end{tcolorbox}}

\newcommand{\kod}[1]{\texttt{\detokenize{#1}}}
\newcommand{\knopka}[1]{\textsf{«#1»}}
\newcommand{\klavisha}[1]{\tcbox[on line, colback=paper, colframe=rule, boxrule=0.4pt, arc=2pt, boxsep=1pt, left=2pt, right=2pt, top=0pt, bottom=0pt]{\sffamily\small #1}}
\newcommand{\galka}{\tikz[baseline=0pt]\draw[ink, line width=0.6pt] (0,0) rectangle (0.27,0.27);}
\usepackage{xspace}
\newcommand{\cm}{\,см\xspace}
\newcommand{\sek}{\,с\xspace}

\hfuzz=1.5pt


\begin{document}

% ---------------------------------------------------------------- титул
\begin{center}
  {\sffamily\small\color{inkgray}Litebee IDE \textbullet{} мастер-класс}\par\vspace{6pt}
  {\sffamily\bfseries\fontsize{26}{30}\selectfont\color{accent}Шоу из четырёх дронов}\par\vspace{6pt}
  {\sffamily\large\color{ink}Пишем программу на Python, проверяем её на экране и отдаём дронам}\par\vspace{8pt}
  \textcolor{accent}{\rule{3.2cm}{1.6pt}}
\end{center}
\vspace{6pt}

Сегодня вы напишете настоящее шоу: четыре дрона взлетят, облетят квадрат, пройдут «волной» по высоте, сделают карусель и сядут каждый на свою клетку. Половина работы уже сделана: в заготовке дрон~1 написан полностью. Ваша задача — по его образцу дописать дронов 2, 3 и 4, а потом, если хватит времени, придумать свою сцену.

\begin{opasno}[Кто запускает дронов]
Летает только ведущий. Ваша задача — шоу, которое проходит проверку в приложении без ошибок и без замечаний о сближении. Готовые шоу ведущий загрузит в дронов и запустит при всех. Пока дроны в воздухе, к коврику не подходим.
\end{opasno}

% ---------------------------------------------------------------- что делаем
\section{Что мы будем программировать}

Шоу длится 54 секунды и состоит из сцен. Все четыре дрона живут по \textbf{одним часам}: сцена начинается у всех в одну и ту же секунду.

\begin{figure}[H]
\centering
\begin{tikzpicture}[x=0.29cm, y=1cm, font=\sffamily\small]
  \foreach \a/\b/\t/\c in {0/5/{калибр.}/paper, 5/11/взлёт/paper, 11/23/{1. Квадрат}/accentsoft, 23/33/{2. Волна}/accentsoft, 33/45/{3. Карусель}/accentsoft, 45/54/посадка/paper}{
    \fill[\c] (\a,0) rectangle (\b,0.8);
    \draw[rule] (\a,0) rectangle (\b,0.8);
    \node[text=ink] at ({(\a+\b)/2},0.4) {\t};
  }
  \foreach \s in {0,5,11,23,33,45,54}{
    \draw[inkgray] (\s,0) -- (\s,-0.12);
    \node[font=\scriptsize, text=inkgray, below] at (\s,-0.1) {\s};
  }
  \node[font=\scriptsize, text=inkgray, anchor=west] at (54.4,-0.32) {с};
\end{tikzpicture}
\caption{План шоу по времени. Голубые сцены вы пишете сами, серые части уже есть в заготовке.}
\end{figure}

\begin{enumerate}
  \item \textbf{Взлёт.} Калибровка, взлёт на 100\cm, каждый дрон зажигает свой цвет и встаёт точно над своей клеткой.
  \item \textbf{Квадрат} (11--23\sek). Каждый дрон облетает квадрат: 4 стороны по 3 секунды и возвращается к себе.
  \item \textbf{Волна} (23--33\sek). Дроны по очереди поднимаются до 140\cm, загораются белым и опускаются обратно.
  \item \textbf{Карусель} (33--45\sek). Все четыре дрона летят по одному кругу вокруг центра квадрата и возвращаются на свои места.
  \item \textbf{Посадка.} Каждый садится точно на свою клетку.
\end{enumerate}

% ---------------------------------------------------------------- коврик
\section{Коврик и клетки}

Дроны летают над ковриком 12 на 12 клеток, каждая клетка 40 на 40\cm. У каждой клетки свой номер от 0 до 143:
\[
  \text{номер клетки} = \text{столбец} + 12 \times \text{строка}.
\]
Например, клетка 52 стоит в столбце 4 и строке 4: $4 + 12 \times 4 = 52$. Мы работаем в одной четверти коврика: столбцы и строки от 0 до 5.

\begin{figure}[H]
\centering
\begin{tikzpicture}[x=1.25cm, y=1.25cm, font=\sffamily]
  \foreach \r in {0,...,5}{
    \foreach \c in {0,...,5}{
      \pgfmathtruncatemacro{\n}{\c + 12*\r}
      \node[font=\scriptsize, text=inkgray, anchor=north west] at (\c+0.03,\r+0.97) {\n};
    }
  }
  \draw[rule] (0,0) grid (6,6);
  \draw[ink, thick] (0,0) rectangle (6,6);
  \foreach \c in {0,...,5}{ \node[font=\small, text=inkgray] at (\c+0.5,-0.3) {\c}; }
  \foreach \r in {0,...,5}{ \node[font=\small, text=inkgray] at (-0.3,\r+0.5) {\r}; }
  \node[font=\small, text=inkgray] at (3,-0.75) {столбец};
  \node[font=\small, text=inkgray, rotate=90] at (-0.75,3) {строка};
  \coordinate (a) at (1.5,1.5); \coordinate (b) at (4.5,1.5);
  \coordinate (c) at (4.5,4.5); \coordinate (e) at (1.5,4.5);
  \draw[accent, very thick, -{Stealth[length=3mm]}, shorten <=9pt, shorten >=9pt] (a) -- (b);
  \draw[accent, very thick, -{Stealth[length=3mm]}, shorten <=9pt, shorten >=9pt] (b) -- (c);
  \draw[accent, very thick, -{Stealth[length=3mm]}, shorten <=9pt, shorten >=9pt] (c) -- (e);
  \draw[accent, very thick, -{Stealth[length=3mm]}, shorten <=9pt, shorten >=9pt] (e) -- (a);
  \draw[inkgray, dashed] (3,3) circle ({1.5*sqrt(2)});
  \draw[accentdark, thick] (2.86,2.86) -- (3.14,3.14) (2.86,3.14) -- (3.14,2.86);
  \draw[inkgray] (3.18,3.0) -- (6.25,3.0);
  \node[font=\small, text=inkgray, anchor=west, align=left] at (6.3,3.0) {центр карусели:\\посередине\\между 13 и 52};
  \foreach \p/\t/\col/\n in {a/13/ledgreen/1, b/16/ledblue/2, c/52/ledyellow/3, e/49/ledred/4}{
    \fill[white, draw=\col, line width=2pt] (\p) circle (0.27);
    \node[font=\small\bfseries, text=ink] at (\p) {\n};
  }
\end{tikzpicture}
\caption{Четверть коврика. Кружки — дроны 1--4 на своих клетках, их цвет — цвет светодиода. Стрелки — маршрут сцены «Квадрат», пунктир — круг карусели.}
\end{figure}

\begin{center}
\small
\begin{tabular}{@{}llll@{}}
\toprule
Дрон & Цвет & Стоит на клетке & Маршрут квадрата \\
\midrule
1 & \kod{'green'} зелёный & 13 & 13 → 16 → 52 → 49 → 13 \\
2 & \kod{'blue'} синий & 16 & 16 → 52 → 49 → 13 → 16 \\
3 & \kod{'yellow'} жёлтый & 52 & 52 → 49 → 13 → 16 → 52 \\
4 & \kod{'red'} красный & 49 & 49 → 13 → 16 → 52 → 49 \\
\bottomrule
\end{tabular}
\end{center}

% ---------------------------------------------------------------- программа
\section{Как устроена программа шоу}

Шоу — это файл на Python. В начале создаётся само шоу, дальше у каждого дрона своя программа: список шагов, которые он делает по порядку, один за другим.

\begin{kod*}
from litebee import Show

show = Show('Мастер-класс: четыре дрона')

d = show.drone(1)      # дальше идут шаги дрона 1
d.calibrate(5)         # калибровка, 5 секунд
d.takeoff(100, 5)      # взлёт на 100 см за 5 секунд
d.rgb('green')         # зажечь зелёный
d.code(13, 1)          # встать над клеткой 13 за 1 секунду
d.code(16, 3)          # перелететь к клетке 16 за 3 секунды
d.land_on(13)          # сесть точно на клетку 13

d = show.drone(2)      # а это уже дрон 2
...
\end{kod*}

\begin{zametka}[Три правила записи]
\begin{itemize}
  \item Каждый шаг начинается с \kod{d.} и заканчивается скобками: \kod{d.code(16, 3)}.
  \item Числа внутри скобок разделяются запятой, а дробная часть пишется через точку: \kod{2.5}, а не \kod{2,5}.
  \item Названия цветов — в кавычках: \kod{d.rgb('red')}. Всё после знака \texttt{\#} — комментарий, дрон его не читает.
\end{itemize}
\end{zametka}

\subsection*{Одни часы на всех}

Почти у каждого шага последнее число — \textbf{сколько секунд он длится}. Шаги идут друг за другом, поэтому время в шоу — это сумма секунд всех шагов до этого места. Если у дрона 1 сцена «Квадрат» начинается на 11-й секунде, то у дронов 2, 3 и 4 тоже должна начинаться на 11-й. Поэтому:

\begin{itemize}
  \item одинаковые сцены у всех дронов должны занимать \textbf{одинаковое число секунд};
  \item если дрону в сцене нечего делать, он \textbf{ждёт}: \kod{d.set_z(100, 4)} — «остаться на высоте 100\cm еще 4 секунды»;
  \item в конце все дроны должны получиться одной длины. Приложение показывает длину каждого дрона в партитуре: полоски должны кончаться вместе.
\end{itemize}

Шаг цвета \kod{rgb} времени не занимает: цвет меняется в этот момент и горит, пока его не сменит следующий \kod{rgb}.

% ---------------------------------------------------------------- команды
\section{Основные команды}

\begin{center}
\small
\begin{tabularx}{\textwidth}{@{}>{\raggedright\arraybackslash}p{4.3cm}X@{}}
\toprule
Команда & Что делает \\
\midrule
\kod{d.calibrate(5)} & Калибровка перед взлётом. Всегда первый шаг. \\
\kod{d.takeoff(100, 5)} & Взлёт на высоту 100\cm за 5 секунд. \\
\kod{d.code(16, 3)} & Лететь по прямой к клетке 16 за 3 секунды. Высота не меняется. \\
\kod{d.set_z(140, 2)} & Подняться или опуститься на высоту 140\cm за 2 секунды. Если высота та же, это \textbf{ожидание} на месте. \\
\kod{d.code_z(27, 130, 4)} & Лететь к клетке 27 и одновременно сменить высоту на 130\cm, за 4 секунды. \\
\kod{d.rgb('blue')} & Зажечь цвет. Есть \kod{'red'}, \kod{'green'}, \kod{'blue'}, \kod{'yellow'}, \kod{'cyan'} (голубой), \kod{'purple'} (фиолетовый), \kod{'white'} и \kod{'black'} (погасить). \\
\kod{d.rgb_value(255, 128, 0)} & Свой цвет: яркость красного, зелёного и синего от 0 до 255. Здесь оранжевый. \\
\kod{d.circle(1, 13, 52, 6)} & Полкруга за 6 секунд. Центр круга — посередине между клетками 13 и 52, радиус — от дрона до центра. Два таких шага подряд — полный круг. \\
\kod{d.land_on(13)} & Сесть точно на клетку 13: подлететь, успокоиться, опуститься и сесть (9 секунд). \\
\bottomrule
\end{tabularx}
\end{center}

\subsection*{Ограничения: за ними следит приложение}

\begin{itemize}
  \item \textbf{Высота} не больше 150\cm.
  \item \textbf{Скорость} не больше 100\cm в секунду. Клетка — 40\cm, значит на одну клетку пути нужно хотя бы 0,4 секунды. Спокойно и красиво — около 1 секунды на клетку.
  \item \textbf{Расстояние} между дронами не меньше 60\cm в любой момент шоу.
  \item \textbf{Ожидание} на месте не дольше 60 секунд за один шаг.
\end{itemize}

% ---------------------------------------------------------------- приложение
\section{Как работать в приложении}

\begin{enumerate}
  \item Откройте Litebee IDE и на стартовом экране нажмите \knopka{Загрузить файл с компьютера}. Выберите файл \kod{zagotovka.py}.
  \item Выберите дрона: нажмите его номер в партитуре или его самого на карте. На вкладке \knopka{Код} откроется его программа. Здесь вы и пишете.
  \item Нажмите \knopka{Проверить}. Приложение выполнит шоу на экране и покажет замечания на вкладке \knopka{Проверки}. Щелчок по замечанию переносит к нужной строке.
  \item Нажмите \knopka{Играть} и смотрите, как летят дроны. На вкладке \knopka{Партитура} видно время: у каждого дрона своя полоска цвета. Перетаскивайте бегунок, чтобы посмотреть любой момент.
\end{enumerate}

\begin{zametka}[Подсказки редактора]
Наберите \kod{d.} — редактор предложит команды. Внутри \kod{d.rgb(} появится палитра цветов: щелчок по цвету сам впишет его в код. Если что-то испортили, \klavisha{Ctrl}+\klavisha{Z} отменяет последнее изменение.
\end{zametka}

\begin{vnimanie}[Пока шоу не дописано]
В заготовке дроны 2--4 ещё не облетают квадрат, поэтому проверка покажет, что дроны сближаются. Это нормально: замечания исчезнут, когда вы допишете сцены.
\end{vnimanie}

% ---------------------------------------------------------------- задания
\section{Задания}

\begin{zadanie}{Квадрат}
Посмотрите, как сцена «Квадрат» написана у дрона 1: четыре шага \kod{code}, по 3 секунды каждый. Допишите её дронам 2, 3 и 4 по таблице маршрутов из раздела «Коврик и клетки». У дрона 2 первая сторона уже написана.

\textbf{Проверка себя.} Нажмите \knopka{Играть}: на 11-й секунде все четыре дрона одновременно трогаются, бегут по квадрату друг за другом и на 23-й секунде стоят на своих местах. Замечаний о сближении больше нет.
\end{zadanie}

\begin{zadanie}{Волна}
Дроны по очереди поднимаются до 140\cm: сначала первый, через 2 секунды второй и так далее. Каждый тратит на подъём 2 секунды и на спуск 2 секунды, остальное время ждёт на высоте 100\cm. Вся сцена занимает \textbf{10 секунд у каждого}.

\begin{center}
\begin{tikzpicture}[x=0.95cm, y=0.62cm, font=\sffamily\small]
  \foreach \n/\w/\col in {1/0/ledgreen, 2/2/ledblue, 3/4/ledyellow, 4/6/ledred}{
    \pgfmathsetmacro{\y}{4-\n}
    \node[anchor=east, text=ink] at (-0.2,\y+0.3) {дрон \n};
    \fill[paper] (0,\y) rectangle (10,\y+0.6);
    \draw[rule] (0,\y) rectangle (10,\y+0.6);
    \fill[\col!55] (\w,\y) rectangle (\w+2,\y+0.6);
    \fill[\col!30] (\w+2,\y) rectangle (\w+4,\y+0.6);
    \node[font=\scriptsize, text=ink] at (\w+1,\y+0.3) {вверх};
    \node[font=\scriptsize, text=ink] at (\w+3,\y+0.3) {вниз};
  }
  \foreach \s in {0,2,...,10}{
    \node[font=\scriptsize, text=inkgray] at (\s,-0.35) {\the\numexpr23+\s\relax};
  }
  \node[font=\scriptsize, text=inkgray, anchor=west] at (10.3,-0.35) {с};
\end{tikzpicture}
\end{center}

Серое — ожидание (\kod{set_z(100, ...)}). У дрона 1 ожидание 6 секунд стоит после подъёма, у дрона 4 — 6 секунд перед подъёмом. Посчитайте ожидания для дронов 2 и 3 сами: до подъёма и после него вместе должно быть 6 секунд. Во время подъёма дрон горит белым, потом возвращает свой цвет.

\textbf{Проверка себя.} В партитуре белые отрезки идут лесенкой, а сцена «Карусель» у всех начинается на 33-й секунде.
\end{zadanie}

\begin{zadanie}{Карусель}
У дрона 1 карусель — это два шага \kod{d.circle(1, 13, 52, 6)}. Всем дронам пишите \textbf{те же самые числа 13 и 52}: это не куда лететь, а две клетки, посередине между которыми находится центр круга. Каждый дрон стоит в своём углу, поэтому сам окажется на своём месте круга.

\textbf{Проверка себя.} На экране четыре дрона идут по одному кругу на одинаковом расстоянии друг от друга и на 45-й секунде снова стоят на своих клетках.
\end{zadanie}

\begin{zadanie}{Посадка и последняя проверка}
Посадка в заготовке уже есть. Убедитесь, что у всех дронов шоу длится 54 секунды (полоски в партитуре кончаются вместе), и пройдите список ниже.

\smallskip
\begin{tabular}{@{}l@{\enskip}l@{}}
\galka & \knopka{Проверить} не показывает ни одной ошибки. \\
\galka & Нет замечаний о том, что дроны сближаются. \\
\galka & Нет замечаний о скорости и высоте. \\
\galka & Все дроны садятся на ту клетку, с которой взлетали. \\
\end{tabular}
\end{zadanie}

\begin{zadanie}{Своя сцена (если осталось время)}
Придумайте четвёртую сцену и вставьте её между каруселью и посадкой, у всех дронов одинаковой длины. Несколько идей:
\begin{itemize}
  \item \textbf{Светофор}: дроны стоят на месте и меняют цвета по очереди (не забудьте ожидание между цветами, иначе цвет не успеет показаться).
  \item \textbf{Лестница}: дроны встают на разные высоты — 80, 100, 120 и 140\cm — с помощью \kod{set_z} или \kod{code_z}.
  \item \textbf{Обратный квадрат}: облететь квадрат в другую сторону.
  \item \textbf{Свои цвета}: оранжевый, розовый, тусклый белый через \kod{rgb_value}.
\end{itemize}
Сцена должна пройти те же проверки, что и всё шоу.
\end{zadanie}

% ---------------------------------------------------------------- ошибки
\section{Если что-то пошло не так}

\begin{center}
\small
\renewcommand{\arraystretch}{1.25}
\begin{tabularx}{\textwidth}{@{}>{\raggedright\arraybackslash}p{6.2cm}X@{}}
\toprule
Приложение пишет & Что сделать \\
\midrule
\kod{Синтаксическая ошибка в строке 42: '(' was never closed} & Не закрыта скобка. Посмотрите на эту строку и строку выше. \\
\kod{NameError: name 'e' is not defined} & Опечатка в имени: \kod{e.code(...)} вместо \kod{d.code(...)}. \\
\kod{Неизвестный цвет 'orange'} & Такого имени цвета нет. Возьмите имя из таблицы команд или \kod{rgb_value}. \\
\kod{Дроны 1 и 2 сближаются до 0 см на 14.0 с} & Два дрона в одно время в одной точке. Обычно один дрон отстаёт по времени или перепутан номер клетки. Сравните его шаги с дроном 1. \\
\kod{... быстрее порога 100 см/с} & Дайте шагу больше секунд. \\
\kod{высота 180 см больше лимита 150 см} & Опустите высоту до 150\cm или ниже. \\
\bottomrule
\end{tabularx}
\end{center}

Проверка не заметит, если вы \textbf{перепутали порядок углов} или \textbf{забыли ожидание} — дроны просто полетят не так, как задумано. Смотрите шоу кнопкой \knopka{Играть} и сравнивайте с дроном 1.

\begin{zametka}[Заберите шоу с собой]
Сохраните файл кнопкой сохранения на верхней панели и заберите его домой. Файл \kod{.py} можно снова открыть в Litebee IDE и продолжить.
\end{zametka}

\end{document}
